% ===================== Chapter 11 =====================
\chapter{Interfaces and Polymorphism}
\label{ch:interfaces}

A \code{Circle} and a \code{Rect} are different structs, with different fields,
yet both have an \emph{area}. Often you want to write code that works with
``anything that has an area'' without caring exactly which shape it is. An
\textbf{interface} captures that idea: it describes a set of methods a type must
provide, and lets you write code against the interface rather than any one
concrete type. This ability to treat many types through one shared description is
called \textbf{polymorphism}.

\section{Defining an interface}

An interface is declared with the \code{interface} keyword and a list of method
\emph{signatures} method names with their parameters and return types, but no
bodies:

\begin{salam}
interface Shape:
    func area() f64
    func name() str
end
\end{salam}

This says: ``a \code{Shape} is anything that has an \code{area()} returning an
\code{f64} and a \code{name()} returning a \code{str}.'' The interface promises
what methods exist; it does not say how they work that is each type's job.

\section{Implementing an interface}

Here is the pleasant part: in Salam a struct implements an interface simply by
\emph{having the right methods}. There is no \code{implements} clause to write and
nothing to declare if a struct has an \code{area()} and a \code{name()} with
matching signatures, it already \emph{is} a \code{Shape}. This is called
\textbf{structural conformance}.

\begin{salam}
interface Shape:
    func area() f64
    func name() str
end

struct Circle:
    r f64 = 0.0
    func area() f64: ret 3.14159 * this.r * this.r end
    func name() str: ret "circle" end
end

struct Rect:
    w f64 = 0.0
    h f64 = 0.0
    func area() f64: ret this.w * this.h end
    func name() str: ret "rect" end
end
\end{salam}

Both \code{Circle} and \code{Rect} provide \code{area()} and \code{name()}, so
both conform to \code{Shape} automatically, without either struct mentioning
the interface. Now we can write code that works with any \code{Shape}.

\section{Static polymorphism: interface bounds}

The first and most efficient way to use an interface is as a \textbf{bound} on a
generic function. The notation \code{<T: Shape>} after a function name means ``for
any type \code{T} that conforms to \code{Shape}.'' Inside the function you may
call any of the interface's methods on a value of type \code{T}:

\begin{salam}
func describe<T: Shape>(s T):
    println s.name(), "area =", s.area()
end

func main:
    describe(Circle { r = 2.0 })
    describe(Rect { w = 3.0, h = 4.0 })
end
\end{salam}

\begin{progout}
circle area = 12.5664
rect area = 12
\end{progout}

The one \code{describe} function serves both shapes. You did not have to write it
twice, and you did not have to tell it which type you are passing the compiler
works that out from the argument. Notice there are no angle brackets at the call
site: you just write \code{describe(Circle \{ r = 2.0 \})}, and Salam infers that
\code{T} is \code{Circle}.

\begin{deep}[In Depth: monomorphization, or ``zero-cost'' generics]
When you call \code{describe(Circle \{...\})}, the compiler generates a
specialised copy of \code{describe} in which \code{T} is exactly \code{Circle},
and another copy for \code{Rect}. Each copy calls the methods directly, with no
run-time lookup it is as fast as if you had written two separate functions by
hand. This technique, called \emph{monomorphization}, is why interface bounds are
described as ``zero-cost'': you get the convenience of writing one generic
function and the speed of hand-specialised code. The price is a little extra
compiled size, one copy per type you actually use.
\end{deep}

\section{Dynamic polymorphism: \texttt{dyn}}

Interface bounds are resolved at compile time, which means each call site knows
the concrete type. But sometimes you want a \emph{single} collection holding a
mixture of shapes some circles, some rectangles and to call \code{area()}
on each without knowing which is which until the program runs. For that, Salam
offers \code{dyn}.

A value of type \code{dyn Shape} is ``some shape, decided at run time.'' You can
store different concrete types in one place and call interface methods on them,
and Salam dispatches each call to the right type's method dynamically:

\begin{salam}
func describe(s dyn Shape):
    println s.name(), "area =", s.area()
end

func main:
    describe(Circle { r = 1.0 })
    describe(Rect { w = 2.0, h = 3.0 })
end
\end{salam}

\begin{progout}
circle area = 3.14159
rect area = 6
\end{progout}

The difference from the bounded version is what happens under the hood: a
\code{dyn Shape} carries, alongside the data, a hidden table of which methods to
call, so the right \code{area()} is chosen at run time. This is exactly what lets
a single list hold a heterogeneous mix of shapes.

\begin{note}
\textbf{When to use which.} Reach for an interface \emph{bound} (\code{<T:
Shape>}) by default it is faster and is the right tool when each call works
with one known type. Reach for \code{dyn} when you genuinely need to mix different
concrete types in one collection or behind one variable, and decide the method at
run time. Bounds are about writing one function for many types; \code{dyn} is
about holding many types in one container.
\end{note}

\section{Designing with interfaces}

Interfaces let you program to a \emph{capability} rather than a concrete type,
which keeps your code flexible. A function that sorts ``anything comparable,'' a
logger that writes to ``anything writable,'' a renderer that draws ``any
shape'' each is expressed by naming the interface it needs and ignoring
everything else about its inputs. When you later add a new type that conforms, the
existing code works with it unchanged, because it only ever depended on the
methods, never the specific struct.

\begin{tip}
Keep interfaces small. An interface with one or two methods (``has an area,''
``can be drawn'') is easy for many types to satisfy and easy to reason about. A
large interface demanding a dozen methods is a burden to implement and tends to
couple your code to one specific design. Prefer several small interfaces over one
big one.
\end{tip}

\section{Chapter summary}

\begin{itemize}[leftmargin=1.4em]
  \item An \code{interface} lists method signatures (no bodies) that a conforming
        type must provide.
  \item Conformance is \emph{structural}: a struct implements an interface just by
        having the matching methods no \code{implements} declaration.
  \item An interface \emph{bound}, \code{func f<T: Iface>(x T)}, writes one
        function usable with any conforming type; the type is inferred at the call
        site and resolved at compile time (monomorphization, zero run-time cost).
  \item \code{dyn Iface} holds ``some conforming type, decided at run time,''
        enabling heterogeneous collections via dynamic dispatch.
  \item Use bounds by default; use \code{dyn} when you must mix concrete types in
        one place.
  \item Favour small interfaces that many types can easily satisfy.
\end{itemize}

\section{Exercises}

\exercise Define an interface \code{Greeter} with a method \code{greet() str}, and
two structs (say \code{Formal} and \code{Casual}) that conform. Write a bounded
function that prints any greeter's greeting.

\exercise Add a third shape, \code{Triangle} (base and height), to the
\code{Shape} example. Confirm the existing \code{describe} works on it with no
changes.

\exercise Give \code{Shape} a third method \code{perimeter() f64} and implement it
for \code{Circle} and \code{Rect}. What error do you get if you forget one?

\exercise Write a bounded function \code{bigger<T: Shape>(a T, b T) str} that
returns the name of whichever shape has the larger area. (Both arguments are the
same type here why?)

\exercise Explain, in a short comment, when you would choose \code{dyn Shape} over
\code{<T: Shape>}, using a list of mixed shapes as your example.
